\section{Smoothing channels on logarithmic and residue space}
\label{sec:channels}

The site model admits a useful smoothing operation: choose a fair subset of
the primes at one site and record the logarithm and a residue of its product.
We prove that this operation suppresses residue dependence in operator norm
and that its logarithmic outputs are uniformly approximable on a fixed
coarse partition.  Operator norm, rather than convergence for each fixed
test, is needed because the single-site tests used later may depend on the
position and on an externally conditioned sample.

All limits in this section are as \(B\to\infty\).  The regularity parameters
fixed in Section~\ref{sec:arithmetic} remain fixed, although no regularity
cutoff is imposed in the channels below.  We use only the prime-product
estimates \eqref{eq:2}--\eqref{eq:3} and the independent site law.
An occurrence of \(x\) inside a channel denotes a logarithmic coordinate,
not the original counting scale.

\subsection{The channel and its two-split kernel}

Let \(\Omega_B\) be the finite set of subsets of \(\mathcal P\), with
probability measure under which the inclusions \(p\in S\) are independent
and have probabilities \(1/p\).  Given \(S\), retain each of its primes
independently with probability \(1/2\), and let \(b\) be their product.
Unconditionally \(b\) has law \(\nu_{1/2}\).
Write
\[
 I=[1/2,3],\qquad x_b=\frac{\log b}{B}.
\]
For a positive integer \(m_1\), partition \(I\) into \(m_1\) equal coarse
intervals \(I_l\).  Choose the number \(n_B\) of fine intervals to be
\[
 n_B=m_1\left\lceil\frac{|I|B^{11/10}}{m_1}\right\rceil,
 \qquad \delta_B=\frac{|I|}{n_B}.
\]
Thus the fine partition refines the coarse partition and
\(\delta_B\asymp B^{-11/10}\) for every fixed \(m_1\).
Use consistent half-open conventions, assigning the final endpoint to
the last cell.

For an integer \(1\le d\le B\), let
\(\mathcal R_d=(\Z/d\Z)^\times\) with uniform probability measure;
for \(d=1\) this is the one-point group.  All products in question are
units modulo \(d\), since \(P_0>B\) for sufficiently large \(B\).
Give \(I\times\mathcal R_d\) the product measure
\[
 d\lambda_d(x,r)=dx\,d{\rm unif}_{\mathcal R_d}(r).
\]
For \(g\in L^2(\Omega_B)\), define a function constant on each fine
logarithmic cell by
\begin{equation}\label{eq:30}
 U_dg(x,r)=\frac{\varphi(d)}{\delta_B}
   \E\!\left[g(S)\ind_{\{x_b\in I_x^{\rm fine},\ b\equiv r\pmod d\}}\right],
 \qquad x\in I_x^{\rm fine}\subset I .
\end{equation}
The expectation in this definition includes both the site and its fair
split.  The domain norm of \(U_d\) is the site \(L^2\) norm, and its range
norm is that of \(L^2(\lambda_d)\).

Let \(P_{\rm av}\) average the residue coordinate.  Then
\[
 u_g:=P_{\rm av}U_dg=U_1g
\]
is independent of \(d\).  Let \(P_{m_1}\) average on the coarse logarithmic
intervals and act identically on the residue coordinate when one is
present.  Let \(P_{\rm fine}\) denote averaging on the fine logarithmic
intervals.  These averaging operators are orthogonal projections, and
\(P_{\rm fine}\) commutes with \(P_{m_1}\).

\begin{proposition}[Uniform channel bounds]\label{prop:channels-operators}
For every fixed coarse partition, uniformly for \(1\le d\le B\),
\begin{equation}\label{eq:31}
 \|U_d\|\ll 1,\qquad
 \|(1-P_{\rm av})U_d\|\ll B^{-c_7},
 \qquad c_7=\frac1{200}.
\end{equation}
Moreover, for every \(\epsilon_1>0\), one can choose \(m_1\) such that
\begin{equation}\label{eq:32}
 \limsup_{B\to\infty}\ 
 \sup_{\|g\|_2\le1}\|u_g-P_{m_1}u_g\|_2\le\epsilon_1 .
\end{equation}
All constants are independent of the input \(g\).  The sufficiently
large threshold for \(B\) may depend on the chosen fixed partition.
\end{proposition}

We first identify the kernel used to prove the proposition.  A cell
\(A=J\times\{r\}\), with \(J\) a fine logarithmic interval, has measure
\[
 \lambda_d(A)=\frac{\delta_B}{\varphi(d)}=:\mu_d.
\]
Set
\[
 p_A(S)=\Prob_{\rm split}(x_b\in J,\ b\equiv r\pmod d\mid S).
\]
For a step function \(F\) in the range space, the adjoint of the channel is
\[
 U_d^*F(S)=\sum_A p_A(S)F|_A.
\]
Consequently \(U_dU_d^*\) has the step kernel
\begin{equation}\label{eq:channels-joint-kernel}
 K_d(A,A')=\frac{\E[p_A(S)p_{A'}(S)]}{\mu_d^2}.
\end{equation}
This is the density on fine cells of two conditionally independent fair
splits of the same site, restricted to \(I\) in each logarithmic coordinate.
The kernel is nonnegative and symmetric.  Its row integral satisfies
\[
 \sum_{A'}K_d(A,A')\lambda_d(A')
 =\frac{\E[p_A(S)\sum_{A'}p_{A'}(S)]}{\mu_d}
 \le\frac{\Prob_{\nu_{1/2}}(x_b\in J,\ b\equiv r\pmod d)}{\mu_d}.
\]
By \eqref{eq:2}, the numerator is at most
\(C\delta_B/\varphi(d)+O(B^{-80})\).  Indeed its logarithmic interval
is contained in the fixed compact interval \(I\), where the densities
are uniformly bounded.  Since
\(\mu_d^{-1}\ll B^{21/10}\), the row integrals are bounded uniformly.
The column integrals have the same bound.  Schur's inequality gives
\(\|U_dU_d^*\|\ll1\), and therefore proves the first assertion of
\eqref{eq:31}.  This proof already applies to every \(L^2\) input,
including signed or complex inputs.

\begin{lemma}[Independent common and exclusive products]
\label{lem:channels-two-splits}
Let \(b_1,b_2\) be two conditionally independent fair splits of the same
site.  Their joint law can be coupled, with failure probability
\(O(P_0^{-1})\), to
\[
 (CE_1,CE_2),
\]
where \(C,E_1,E_2\) are independent products with law \(\nu_{1/4}\).
Replacing the two-split step kernel by this independent-product step
kernel changes its operator norm by \(O(B^5/P_0)\), uniformly for \(d\le B\).
\end{lemma}

\begin{proof}
At a fixed prime \(p\), the categories \emph{common to both splits},
\emph{exclusive to the first}, and \emph{exclusive to the second} have
probabilities \(1/(4p)\) each, and are mutually exclusive.
The joint law of three independent Bernoulli indicators with these
marginal probabilities differs from this categorical law in total
variation by \(O(p^{-2})\).  For example, the probability of two or more
independent successes is \(O(p^{-2})\), and each remaining probability
differs from its categorical value by \(O(p^{-2})\).
Couple these laws independently over the primes.  A union bound gives
total failure probability
\[
 O\!\left(\sum_{p>P_0}p^{-2}\right)=O(P_0^{-1}).
\]
On successful coupling the products agree.  Collisions that would put
a squared prime in one of the independent products are included in the
failure event.

If two joint probability laws differ in total variation by \(\varepsilon\),
each cell-pair probability differs by at most a constant times
\(\varepsilon\).  Their step kernel difference is thus bounded pointwise
by \(O(\varepsilon\mu_d^{-2})\).
The total measure of the output space is \(|I|\), so Schur's inequality
bounds the operator difference by the same expression times \(|I|\).
Now
\[
 \mu_d^{-2}\ll B^{42/10}\ll B^5.
\]
Taking \(\varepsilon=O(P_0^{-1})\) proves the assertion.
\end{proof}

\subsection{Removing small exclusive products}

The two-split model reduces the channel operator to independent common
and exclusive products. An exclusive product near $1$ provides too little
averaging to erase its residue, so we first bound the operator contribution
of those products. Both row and column bounds are needed.

For a product \(v\), write \(\ell_B(v)=(\log v)/B\).
Let \(K_d^{\rm ind}\) be the independent-product kernel from
Lemma~\ref{lem:channels-two-splits}.
Its entries are
\[
 K_d^{\rm ind}(A,A')
 =\frac{\Prob((\ell_B(CE_1),CE_1\bmod d)\in A,\
              (\ell_B(CE_2),CE_2\bmod d)\in A')}{\mu_d^2}.
\]
The following estimate concerns the operator obtained by retaining only
the specified event in this probability.

\begin{lemma}[Two-sided Schur estimate for small exclusives]
\label{lem:channels-low-exclusive}
For \(0<\theta\le1/10\), the contribution to \(K_d^{\rm ind}\) from
\(\ell_B(E_i)\le\theta\), for either \(i=1\) or \(i=2\), has operator norm
\[
 \ll (\theta+1/R)^{1/4}+B^{-70}+B^5/P_0,
\]
uniformly for \(d\le B\).  The same bound, with a different absolute
constant, holds for the union of the two events.
\end{lemma}

\begin{proof}
It suffices to consider \(i=2\).  Put
\[
 \sigma_\theta=\Prob_{\nu_{1/4}}(\ell_B(E_2)\le\theta)
 \ll(\theta+1/R)^{1/4},
\]
where the last inequality is \eqref{eq:3}.
The row integral at \(A=J\times\{r\}\) of the restricted kernel is at most
\[
 \frac{\Prob((\ell_B(CE_1),CE_1\bmod d)\in A,\
                         \ell_B(E_2)\le\theta)}{\mu_d}
 =\sigma_\theta\,
   \frac{\Prob((\ell_B(CE_1),CE_1\bmod d)\in A)}{\mu_d}.
\]
Here \(E_2\) is independent of \(C,E_1\).
The marginal law of \(CE_1\) differs from \(\nu_{1/2}\) by
\(O(P_0^{-1})\) in total variation: the independent common and exclusive
indicators can be coupled to their mutually exclusive counterparts
prime by prime exactly as in the preceding lemma.  Thus \eqref{eq:2}
and \(\mu_d^{-1}\ll B^{21/10}\) bound the last display by
\[
 C\sigma_\theta+O(B^5/P_0)+O(B^{-70}).
\]

For the column integral at \(A'=J'\times\{r'\}\), first omit the condition
that \(CE_1\) has logarithm in \(I\).  Condition on \(E_2=e\) with
\(\ell_B(e)\le\theta\).  The remaining condition on \(C\) is
\[
 \ell_B(C)\in J'-\ell_B(e),\qquad
 C\equiv r'e^{-1}\pmod d.
\]
The shifted interval lies in \([2/5,3]\), has length \(\delta_B\),
and the prescribed residue is a unit.  Applying \eqref{eq:2} with
\(z=1/4\), uniformly in \(e\), gives conditional probability at most
\[
 C\frac{\delta_B}{\varphi(d)}+O(B^{-80}).
\]
Integration over the event \(\ell_B(E_2)\le\theta\) and division by
\(\mu_d\) therefore give a column bound \(C\sigma_\theta+O(B^{-70})\).
Schur's inequality proves the claimed operator bound.
The \(i=1\) case follows on transposing the kernel.  The kernel of the
union is nonnegative and is entrywise at most the sum of the two
restricted kernels, so the row and column estimates also prove its bound.
\end{proof}

The separate column estimate is essential: an estimate for the total
probability of a discarded event would not by itself control its
operator norm.

\subsection{Flattening all nonconstant residue modes}

We prove the second assertion of \eqref{eq:31}.
Set \(\theta=B^{-1/10}\) and retain the part of \(K_d^{\rm ind}\) where
\(\ell_B(E_1),\ell_B(E_2)>\theta\).
For a fixed common product \(C=c\), the probability that one retained
exclusive lands in the output cell \(A=J\times\{r\}\) is
\[
 p_A^{(\theta)}(c)
 =\Prob\bigl(\ell_B(E)\in(J-\ell_B(c))\cap(\theta,\infty),\
                   E\equiv rc^{-1}\pmod d\bigr).
\]
If the interval is nonempty it is contained in
\([B^{-1/10},3]\), since \(\ell_B(c)\ge0\).  The interval and its endpoint
conventions, including a partial cell at \(\theta\), are within the
uniform statement of \eqref{eq:2}.  Consequently
\begin{equation}\label{eq:channels-exclusive-cell}
 p_A^{(\theta)}(c)
 =\frac1{\varphi(d)}
   \int_{(J-\ell_B(c))\cap(\theta,\infty)}f_{1/4,B}(t)\,dt
   +O(B^{-80}).
\end{equation}
The main term depends on \(J\) and \(\ell_B(c)\), but not on \(r\) or
on the residue of \(c\).  All terms are uniformly bounded: the
probability is at most one and its main term differs by \(O(B^{-80})\).
Alternatively the bound \(f_{1/4,B}(t)\ll t^{-3/4}\) gives the explicit
cell bound \(O(\delta_B B^{3/40}/\varphi(d))\).

Conditional on \(C\), the exclusives are independent.  Multiplying
\eqref{eq:channels-exclusive-cell} for two output cells and integrating
over \(C\), the retained kernel differs from a kernel independent of
both residues by at most
\[
 O(B^{-80}\mu_d^{-2})=O(B^{-75})
\]
pointwise, and hence by \(O(B^{-70})\) in operator norm.
A kernel independent of both residues is annihilated by projection
onto \(1-P_{\rm av}\) on either side.  Lemmas
\ref{lem:channels-two-splits} and \ref{lem:channels-low-exclusive}
therefore yield
\[
 \begin{split}
 \|(1-P_{\rm av})U_dU_d^*(1-P_{\rm av})\|
 &\ll (B^{-1/10}+1/R)^{1/4}
          +B^{-70}+B^5/P_0\\
 &\ll B^{-1/40},
 \end{split}
\]
because \(1/R=(\log P_0)/B=O(\log B/B)\) and \(P_0=B^{1000}\).
Taking the square root of this operator identity gives the stronger bound
\[
 \|(1-P_{\rm av})U_d\|\ll B^{-1/80}.
\]
In particular \eqref{eq:31} holds with \(c_7=1/200\).
This reasoning did not fix an input function at any stage, so the
estimate is uniform on the whole unit ball of \(L^2(\Omega_B)\).

\subsection{Uniform coarse compactness}

Residue dependence is now negligible uniformly in the input function.
The remaining task is to approximate the logarithmic output uniformly on
a fixed finite partition. This will give a finite family of endpoint
features for the graph kernel.

We now prove \eqref{eq:32}.  It suffices to use \(d=1\), since \(u_g=U_1g\).
Fix \(0<\theta<1/20\) and choose a smooth function \(a_\theta:\R\to[0,1]\)
equal to zero on \((-\infty,\theta]\) and to one on
\([2\theta,\infty)\).  It may be chosen nondecreasing.
In the independent-product kernel insert the weight
\(a_\theta(\ell_B(E_1))a_\theta(\ell_B(E_2))\).
The change is a nonnegative kernel supported where at least one
exclusive logarithm is at most \(2\theta\).  Lemma
\ref{lem:channels-low-exclusive} bounds its operator norm by
\[
 O((2\theta+1/R)^{1/4})+O(B^{-70}+B^5/P_0).
\]

Choose also a smooth upper cutoff equal to one on \([0,3]\) and supported
below \(16/5\).  On the positive axis let \(f_{1/4,B}^{\rm cut}\) be the
product of this upper cutoff, \(a_\theta\), and \(f_{1/4,B}\), extended
by zero to the real line.  For each fixed \(\theta\), all derivatives
of any fixed order of this function are bounded uniformly in \(B\).
This follows from the derivative bounds for $f_{1/4,B}$ established
before Proposition~\ref{prop:arith-local-laws};
the cutoff keeps the argument away from zero.
The upper cutoff changes none of the probabilities relevant to the
output interval \(I\).

For completeness, the weighted version of the local law used here
follows directly from its interval version.  For fixed \(z=\ell_B(c)\)
and a fine interval \(J\), apply \eqref{eq:2} to subintervals of
\((J-z)\cap[\theta,3]\).  The difference between the exclusive-product
measure and the density \(f_{1/4,B}(t)\,dt\) has cumulative integral
\(O(B^{-80})\) there.  Stieltjes integration by parts against
\(a_\theta\) bounds its weighted integral by \(O_\theta(B^{-80})\):
the endpoint values and the total variation of \(a_\theta\) are bounded.
Thus, uniformly in \(c,J\),
\[
 \E\!\left[
   a_\theta(\ell_B(E))\ind_{\{\ell_B(cE)\in J\}}\right]
 =\int_J f_{1/4,B}^{\rm cut}(x-\ell_B(c))\,dx
       +O_\theta(B^{-80}).
\]
Conditional independence of the two exclusives now identifies the
weighted step kernel, up to \(O_\theta(B^{-70})\) in operator norm,
with the fine-cell compression of the continuous kernel
\begin{equation}\label{eq:channels-smooth-kernel}
 K_{B,\theta}(x,y)=
 \E_C\!\left[
  f_{1/4,B}^{\rm cut}(x-\ell_B(C))
  f_{1/4,B}^{\rm cut}(y-\ell_B(C))\right],
 \qquad (x,y)\in I^2.
\end{equation}
Indeed multiplying the two weighted cell formulas makes an
\(O_\theta(B^{-80})\) error in each cell-pair probability; division by
\(\delta_B^2\) still leaves an operator error smaller than
\(O_\theta(B^{-70})\).

The kernel in \eqref{eq:channels-smooth-kernel} and its first derivatives
are bounded uniformly in \(B\), with constants depending on \(\theta\).
Differentiation may be taken under the expectation because the
cut densities and their derivatives have uniform bounds and the law
of \(C\) is a probability measure.  In particular no limiting law for
\(C\), and no derivative bound for that law, is needed.
Let \(\mathcal T_{B,\theta}\) be the integral operator with this kernel.
If \(a=|I|/m_1\) is the coarse mesh, the mean value theorem gives
\[
 \sup_{x,y\in I}
 \left|K_{B,\theta}(x,y)
      -(P_{m_1}^{(x)}P_{m_1}^{(y)}K_{B,\theta})(x,y)\right|
 \ll_\theta a.
\]
Here the two superscripts denote averaging the indicated kernel
coordinate.  Schur's inequality consequently gives
\begin{equation}\label{eq:channels-coarse-kernel}
 \|\mathcal T_{B,\theta}
           -P_{m_1}\mathcal T_{B,\theta}P_{m_1}\|
 \ll_\theta m_1^{-1}.
\end{equation}

Write \(Q=1-P_{m_1}\).  Fine and coarse averaging commute, and fine
averaging is a contraction.  Therefore
\[
 \|Q P_{\rm fine}\mathcal T_{B,\theta}P_{\rm fine}Q\|
 \le \|Q\mathcal T_{B,\theta}Q\|
 \ll_\theta m_1^{-1}.
\]
Combining this with the independent-model comparison and the cutoff
estimate proves, for every fixed \(\theta,m_1\),
\[
 \begin{split}
 \|Q U_1\|^2
 &=\|Q U_1U_1^*Q\|\\
 &\le C(2\theta+1/R)^{1/4}
        +C_\theta m_1^{-1}+o_{\theta,m_1}(1).
 \end{split}
\]
Given \(\epsilon_1>0\), first choose \(\theta\) so that
\(C(2\theta)^{1/4}<\epsilon_1^2/4\), and then choose \(m_1\) so that
\(C_\theta/m_1<\epsilon_1^2/4\).
Taking the upper limit in \(B\) and the square root proves
\eqref{eq:32}, and completes the proof of
Proposition~\ref{prop:channels-operators}.

\begin{corollary}[Coarse site features]\label{cor:channels-features}
For the chosen coarse partition, define
\[
 V_l(S)=\frac1{|I_l|}
          \Prob_{\rm split}(x_b\in I_l\mid S).
\]
Then \(0\le V_l\le |I_l|^{-1}\), and for every \(g\in L^2(\Omega_B)\)
the coarse value of \(P_{m_1}u_g\) on \(I_l\) is exactly
\[
 \E[g(S)V_l(S)].
\]
In particular the bounds \eqref{eq:31}--\eqref{eq:32} and this identity
hold uniformly when \(g\) varies with a position, with \(B\), or with
an external parameter.
\end{corollary}

\begin{proof}
Integrating \eqref{eq:30} with \(d=1\) over \(I_l\) sums precisely the
fine cells contained in \(I_l\).  The result is
\(\E[g(S)\Prob_{\rm split}(x_b\in I_l\mid S)]\).
Divide by \(|I_l|\).  The bounds on \(V_l\) follow from its definition;
the final uniformity follows because the preceding results are
operator bounds rather than fixed-input limits.
\end{proof}
